\section{Results}
\label{sec:results}

We report the four audit stages of \S\ref{sec:method} on discovery data;
aggregation, per-head estimates and sensitivity analyses are detailed in
Appendices~\ref{app:stage1}--\ref{app:stage4}.

\subsection{RI poorly identifies causal contributors}
\label{sec:results-why}

\paragraph{What RI selects.}
Under our percentile rule, pooled RI selects two heads whose qualifying attention
is exclusively self- or previous-token attention; demonstrations contribute
over 99\% of their scores (Table~\ref{tab:pooled}), consistent with static
token preferences rather than retrieval. No eligible head passes the matched-target
test after Holm correction.
The test-only rule selects 59 heads, but several
are sensitive to the current token, target anchor or fact order
(Appendix~\ref{app:stage1}).

\begin{figure}[t]
\centering
\includegraphics[width=\columnwidth]{fig_ri_vs_importance}
\caption{RI versus exact signed Scope-P importance on 40 common pairs
(symmetric-log axis). Orange marks the 61 RI-selected heads; shading marks
$|I_h|<0.3$. (a) Pooled RI for all 1,024 heads. (b) First-token test-only
name gap for 167 supported heads.}
\label{fig:ri-vs-imp}
\end{figure}

\paragraph{Against the causal map.}
Pooled RI correlates weakly with signed importance ($\rho=0.13$), whereas
the test-only name gap correlates negatively ($\rho=-0.25$, or $-0.30$
within layer; Figure~\ref{fig:ri-vs-imp}). The two pooled-index heads have
near-zero effects, and 48 of the combined 61 selected heads have
$|I_h|<0.1$, less than 1\% of the clean--corrupted gap. Conversely, among
the 148 heads measured on all 178 pairs, only eight of the 25 heads with
$|I_h|\geq0.3$ are selected. Both selections are therefore incomplete,
and the pooled selection also meets our pre-registered misleading criterion
(Appendix~\ref{app:stage2}).

\paragraph{Why RI can miss retrieval.}
At the answer position the current token is always a colon, so with visible
token IDs unchanged, raw-embedding RI cannot adapt its name preferences when
mothers are reassigned. Copying a current non-target name may also raise the
name-control score and penalize useful heads, a hypothesis tested in
\S\ref{sec:results-char}. These limitations concern the index, not the
model's capacity for semantic computation.

\subsection{The causal map}
\label{sec:results-map}

\paragraph{A few heads carry most of the effect.}
The gradient screen closely reproduces the exact ranking
($\rho=0.945$; Appendix~\ref{app:stage2}), and all effects reported below use
exact patching. Against a mean clean--corrupted gap of 11.2 logits, replacing
L17H1 changes the metric by 5.7 logits, about half the gap, while three further
heads each change it by roughly a quarter (Figure~\ref{fig:causal-map}).
Twenty-five heads reach $|I_h|\geq0.3$: 17 positive and eight negative, with
negative effects shifting the model toward the stated fact under corrupted-output
insertion. These are not
independent shares because head effects can interact. The same four heads lead
the exact common-subset scan, indicating that the leading heads are not an artefact of
the extension set (Appendix~\ref{app:stage2}).

\begin{figure}[t]
\centering
\includegraphics[width=\columnwidth]{fig_causal_map}
\caption{The causal map: signed Scope-P importance of the 25 strong heads
on 178 pairs, with family-bootstrap intervals; orange marks RI-selected heads.}
\label{fig:causal-map}
\end{figure}

\paragraph{Two positions, two kinds of head.}
Restricting the patch to the final colon splits the map by depth
(Figure~\ref{fig:scopeF-ratio}, Appendix~\ref{app:stage2}). In the Scope-F set,
measurable heads in layers 16--26 retain nearly all their
Scope-P effect under Scope F, indicating answer-position contributions, whereas
heads in layers 6--11 lose their effects and therefore write information at
earlier positions for later components to read; two layer-14--15 heads are
intermediate. At the colon, positive heads raise the donor-retrieved name and
negative heads lower it (Table~\ref{tab:scopeF-app});
\S\ref{sec:results-char} extends localization to all strong heads and tests
how their effects arise.

Together, these positional differences suggest a writer--reader organization
whose computations and connections are tested in
\S\S\ref{sec:results-char}--\ref{sec:results-routes}.

\subsection{What the important heads do}
\label{sec:results-char}
\paragraph{Where the heads act, and what the writer writes.}
Twenty of the 25 strong heads deliver nearly all their effect at the colon
($I^F_h/I_h\in[0.96,1.07]$; Appendix~\ref{app:stage3}). The strongest head,
L17H1, instead acts within the query fact, principally at the mother tokens,
``is'' and the period (Figure~\ref{fig:profiles}). Because some of these
positions have identical clean and corrupted tokens, its output carries
information about the preceding mother. At ``is'', replacing L17H1 shifts a
controlled mother-token readout by $-0.46$, compared with $-0.52$ under a full
source swap, with a negative shift in all 20 families. This supports an
identity-sensitive write. A second writer,
L15H25, acts at the final token of the query child; its content remains unresolved
because even the full source swap barely changes its readout
(Appendix~\ref{app:stage3}). The strongest writer thus shares the readers' layer
band; the shallow heads above are weak.

\begin{figure}[t]
\centering
\includegraphics[width=\columnwidth]{fig_head_profiles}
\caption{Single-position effects of the three fact-position
heads summed by token role (M: mother, C: child, oth.: other test facts,
Q: question), averaged over 20 common families; error bars show SD.}
\label{fig:profiles}
\end{figure}

\paragraph{What the readers select and what they carry.}
Answer-position heads form three attention profiles at the colon
(Figure~\ref{fig:colon-attention}, Appendix~\ref{app:stage3}): mother-attenders,
child-attenders that select the queried child's second mention, and colon-attenders. Under a question-only change,
mother- and child-attenders shift to the newly queried fact in all 20 common
families, whereas reordering or swapping mothers barely changes family-mean
attention. Pattern--value
separation shows that values carry most of the leading heads' mean effects: source
selection follows the question, while the swap changes what is read there.
Output projections explain only part of this computation: L27H6 directly favors
the answer, but several strong positive, negative and child-attending heads have
projections much smaller than their causal effects. Replacing the raw embedding
with each head's contextual input improves the correlation with importance only
from $-0.26$ to $-0.17$ over 68 supported heads and therefore does not repair RI
(Appendix~\ref{app:stage3}).

\paragraph{Copying is part of the story, not the mechanism.}
Under the first-token criterion, six heads promote the name they attend to and one
suppresses it, but all such labels disappear at the last-token anchor. For positive
movers, the RI name gap becomes negative when the current token is a control name,
supporting the copy-penalization hypothesis of
\S\ref{sec:results-why}. Nevertheless, copying is neither necessary nor
sufficient: the strongest weight copier copies the distractor, while L23H15 favors
its attended name despite a causal effect of $-0.95$. Thus projection sign does
not determine causal sign, and ``negative mover'' does not identify a single
mechanism (Appendix~\ref{app:stage3}).

These profiles define the Stage-4 hypotheses: query-fact writers may feed
mother- or child-attending readers, while the colon heads and positive and
negative movers require route and group tests
(\S\ref{sec:results-routes}).
